\documentclass[a4paper,12pt]{article}
\usepackage[margin=.7in]{geometry}

% Encoding and font settings
\usepackage[utf8]{inputenc}

% Math packages
\usepackage{amsmath,amsfonts,amssymb,amsthm}
\usepackage{bbm}
% Formatting packages
\usepackage{indentfirst,color,fullpage}

% Graphics and tables
\usepackage{graphicx}
\usepackage[flushleft]{threeparttable}
\usepackage{booktabs}

% Hyperlinking and referencing
\usepackage{hyperref}
\usepackage{natbib}
\usepackage{url}

% Additional functionality
\usepackage{footnote}
\usepackage{xcolor}
\usepackage{bm}
\usepackage{ulem}
\newcommand{\red}{\textcolor{red}}
\newcommand{\purple}{\textcolor{purple}}

\usepackage{todonotes}

% Define custom maths equations 
\DeclareMathOperator{\KL}{KL}
\DeclareMathOperator{\JS}{JS}
\DeclareMathOperator*{\argmin}{arg\,min}
% \DeclarePairedDelimiter\abs{\lvert}{\rvert}%
\DeclareMathOperator{\EX}{\mathbb{E}}% expected value
\newcommand{\JSGa}{\mathop{\mathrm{JS}^{\mathrm{G}_{\lambda}}}\nolimits}
\newcommand{\cJSGa}[1]{\mathop{\mathrm{cJS}_{i, #1}^{\mathrm{G}_{\lambda}}}\nolimits}
\newcommand{\eJSGa}[1]{\mathop{\mathrm{eJS}_{i, #1}^{\mathrm{G}_{\lambda}}}\nolimits}
\DeclareRobustCommand{\eJSGai}{\mathop{\mathrm{eJS}^{\mathrm{G}_{\lambda}}}\nolimits}
\DeclareMathOperator{\GJS}{\mathrm{JS}^{\mathrm{G}}}
\DeclareMathOperator{\CE}{\mathrm{CE}}
\newcommand{\rhou}[1]{\rho_{\text{#1}}}
\newcommand\numberthis{\addtocounter{equation}{1}\tag{\theequation}}

\newcommand{\mvmean}[1]{\symbf{\alpha}_{\mid i, #1} y + y^{\symbf{\beta}_{\mid i, #1}} \symbf{\mu}_{\mid i, #1}}
\newcommand{\mvcov}[1]{\left(y^{\symbf{\beta}_{\mid i, #1}}\right)^{\mathrm{T}} \Sigma_{\mid i, #1} y^{\symbf{\beta}_{\mid i, #1}}}



\begin{document}

\noindent \textbf{{\Large Response to Reviewers:\\
% ``MOPED: A moving sum method for change point detection in pairwise extremal dependence''}}
``Clustering of multivariate tail dependence using conditional methods''}}

\vspace*{.25cm}

% \noindent \textbf{{\large(Manuscript ID: Extremes b660ab30-cacb-4a14-a6a9-e60fcf48c005)}}
\noindent \textbf{{\large(Manuscript ID: CSDA-D-25-01358)}}

\vspace*{.25cm}

% \noindent {\it by Euan T.~McGonigle, Matthew Pawley, Jordan Richards\footnote{Corresponding author: Jordan Richards \\
\noindent {\it by Patrick O'Toole, Christian Rohrbeck, Jordan Richards\footnote{Corresponding author: Patrick O'Toole \\
\hspace*{.75cm} Department of Mathematical Sciences, University of Bath, UK\\
\hspace*{.75cm} \href{mailto:pot23@bath.ac.uk}{pot23@bath.ac.uk}}}

\vspace{.5cm}

% \noindent {\it To: Rafal Kulik, Ph.D.} \\
% {\it Editor, Extremes} \\
% {\it Chen Zhou, Ph.D.} \\
% {\it Guest Editor
% } \\

% Dear Rafal and Chen, \\
Dear Referees,

% Thank you for the opportunity to provide an updated version of our manuscript entitled ``MOPED: A moving sum method for change point detection in pairwise extremal dependence'' for consideration in Extremes. The comments from the reviewers have been very helpful, and we have addressed them thoroughly.\\
Thank you for the opportunity to provide an updated version of our manuscript entitled ``Clustering of multivariate tail dependence using conditional method'' for consideration in Computational Statistics and Data Analysis. The comments from the reviewers have been very helpful, and we have addressed them thoroughly.\\

Major changes to the text have been highlighted in \textcolor{red}{red}.
\clearpage
\newpage

\clearpage
\newpage
%%%%%%%%%%%%%%%%%%%%%%%%%%%%%%%%%%%%%%%%%%%%%%%%
%%%%%%%%%%%%%% End of Highlights %%%%%%%%%%%%%%%
%%%%%%%%%%%%%%%%%%%%%%%%%%%%%%%%%%%%%%%%%%%%%%%%


%%%%%%%%%%%%%%%%%%%%%%%%%%%%%%%%%%%%
%%% Editor %%% Editor %%% Editor %%%
%%%%%%%%%%%%%%%%%%%%%%%%%%%%%%%%%%%%
\noindent {\Large \bf {Responses to the Co-Editor}}\\

% \noindent \textcolor{violet}{I have received reports from two experts in the field. Referee 1 recommended "Revise" with a list of comments while Referee 2 is more critical.\\
% The main critique of Referee 2  is that the paper is lack of theoretical guarantee. To my view, it is also important to have methodological papers ahead of theoretical work. Therefore, I would not emphasize on developing new theories. Instead, it is sufficient to provide arguments why the proposed procedure has the potential to achieve theoretical guarantee.\\
% With this idea in mind, I'd like to invite the authors for a revision addressing the comments of Referee 1 and the minor comment of Referee 2.}\\
\noindent \textcolor{violet}{Your article was reviewed by two experts in the field. The topic is clearly interesting, but the article requires a major revision taking into account all the reviewers' comments, especially those of Reviewer 1. Please provide a letter describing all your changes and responses to the reviewers' comments.}\\

% \noindent \textcolor{violet}{\bf Response}: Thank you for these comments. We have thoroughly addressed the comments of Reviewer 1, and so the latest version of the manuscript includes two additional simulation studies, a re-write of the methodology section, and additional comments on model assumptions and justification throughout the paper. Per your advice, we have opted to leave the major comments of Reviewer 2 unaddressed, but we have addressed all of the minor comments. \textbf{Note that we have not derived any new theoretical results in this updated manuscript.}\\

\todo{Include what we plan to do here!}
\noindent \textcolor{violet}{\bf Response}: Thank you for these comments. We have thoroughly addressed the comments of $\ldots$ \\

\clearpage
\newpage

%%%%%%%%%%%%%%%%%%%%%%%%%%%%%%%%%%%%%%%%%%%%%%%%
%%% Reviewer 1 %%% Reviewer 1 %%% Reviewer 1 %%%
%%%%%%%%%%%%%%%%%%%%%%%%%%%%%%%%%%%%%%%%%%%%%%%%
\noindent {\Large {\bf Responses to Reviewer 1}}\\

% \noindent \textcolor{teal}{The paper is well-written, interesting, and potentially useful. However, several assumptions are not stated explicitly, some imported results are used
% without verifying their conditions in the present setting, and the theoretical
% guarantees (e.g., consistency for break detection) are not yet articulated.
% Clarifying these points and expanding the Monte Carlo section would substantially
% strengthen the paper.}
\noindent \textcolor{teal}{
While a substantial body of work exists on the estimation of conditional extremes models in multivariate and spatial contexts, the problem of clustering within this framework has received little attention. 
The paper under review aims to address this gap by developing clustering algorithms for the analysis of conditional extremes models
across a broad range of settings, with the goal of supporting exploratory analysis and parameter estimation. 
While the paper addresses an important and innovative topic, the current version suffers from significant theoretical, methodological, and empirical limitations. 
The use of heterogeneous data sources without proper justification or harmonization, the lack of validation of key modeling assumptions, and the unclear presentation of the dissimilarity measure and simulation settings are major concerns. 
In a subsequent revision, my opinion is that the authors must:
\begin{itemize}
  \item Clarify and justify the properties and interpretability of the proposed dissimilarity measure.
  \item Revise the simulation study to clearly demonstrate the method’s robustness and relevance within the CEV framework.
  \item Provide rigorous validation of the Gaussian assumption for conditional residuals in the real data application, and either justify the use of Met Éireann and ERA5 data or switch to a single, consistent data source.
  \item Address all minor points regarding clarity and presentation.
\end{itemize}
The acceptability of the revised manuscript will depend heavily on the authors’ ability to address these concerns satisfactorily. 
This is a critical point: the scientific validity of the results hinges on
these aspects.
}

%%% Section %%% Section %%% Section %%%
%%% Section %%% Section %%% Section %%%
%%% Section %%% Section %%% Section %%%
\section*{{Major Comments}}

%%% Item %%% Item %%% Item %%% Item %%% Item %%% Item %%% Item %%%
% \noindent \textcolor{teal}{1. The way you introduce the multivariate time series (viz., Eq. (4)) reads as if a break affects the entire joint distribution, not just extremal
% dependence. Please clarify.}\\
% 
% \noindent \textcolor{teal}{\bf Response}: See below. \\

\noindent \textcolor{teal}{A first concern is the limited readability of Section 2 and the fact that the proposed dissimilarity
measure is not sufficiently discussed in terms of its properties and interpretability. 
Further details are provided below.:}

% \noindent \textcolor{teal}{
%   1. The statement on page 9 claiming that $\mathrm{G}_{\lambda}(H, H^*)$ is multivariate Gaussian is unclear and potentially misleading. 
%   As currently phrased, it may be interpreted as follows: let $X \sim \mathcal{N}_d(\mu_1, \Sigma_1)$ and $Y \sim \mathcal{N}_d(\mu_2, \Sigma_2)$ be two multivariate Gaussian random vectors with cumulative distribution functions $H$ and $H^*$, respectively. 
%   Then the weighted geometric mean
% \[
%   \mathrm{G}_{\lambda}(H, H^*) = H^{1-\lambda}(H^*)^{\lambda}
% \]
% would itself correspond to a multivariate Gaussian distribution. 
% This interpretation is incorrect in general. 
% Presumably, the authors intend to refer instead to the normalized weighted geometric mean of the densities, namely
% \[
%   \frac{f^{\lambda}h^{1-\lambda}}{\int f^{\lambda}h^{1-\lambda} dx},
% \]
% which is indeed Gaussian. 
% Clarifying this point would help avoid confusion.
% }

{\color{teal}
\noindent 1. The statement on page 9 claiming that $\mathrm{G}_{\lambda}(H, H^*)$ is multivariate Gaussian is unclear and potentially misleading. 
As currently phrased, it may be interpreted as follows: let $X \sim \mathcal{N}_d(\mu_1, \Sigma_1)$ and $Y \sim \mathcal{N}_d(\mu_2, \Sigma_2)$ be two multivariate Gaussian random vectors with cumulative distribution functions $H$ and $H^*$, respectively. 
Then the weighted geometric mean
\[
  \mathrm{G}_{\lambda}(H, H^*) = H^{1-\lambda}(H^*)^{\lambda}
\]
would itself correspond to a multivariate Gaussian distribution. 
This interpretation is incorrect in general. 
Presumably, the authors intend to refer instead to the normalized weighted geometric mean of the densities, namely
\[
  \frac{f^{\lambda}h^{1-\lambda}}{\int f^{\lambda}h^{1-\lambda}\,dx},
\]
which is indeed Gaussian. 
Clarifying this point would help avoid confusion.
}

\noindent \textcolor{teal}{\bf Response}: \\

%%% Item %%% Item %%% Item %%% Item %%% Item %%% Item %%% Item %%%

{\color{teal}
\noindent 2. Could the authors discuss possible advantages or interpretations of the expected divergence defined in Equation (11)? 
At the beginning of Section 2.3, four desirable properties of a divergence measure are listed: non-negativity, symmetry, being upper-bounded, and computational efficiency. 
Only the second property (symmetry) is briefly addressed. 
Do the other three properties hold for the proposed divergence? 
Regarding interpretability, what can be inferred about $Y_{i,s}$ and $Y_{i,s^*}$ if
$\eJSGa{i, s, s^*} = 0$?
Do the two variables then have identical conditional $(d - 1)$-variate Gaussian distributions at
sites $s$ and $s^*$, respectively? 
If so, do we also have equivalence?
}

\noindent \textcolor{teal}{\bf Response}: \\

%%% Item %%% Item %%% Item %%% Item %%% Item %%% Item %%% Item %%%

\noindent \textcolor{teal}{My second concern relates to the simulation study, for which it is unclear how the considered setup fits within the framework introduced in Section 2.}

\todo{Fix references here?}
\noindent \textcolor{teal}{1. In Section 3.1, in the Gaussian copula setting, Keef, Papastathopoulos, and Tawn 2013, Section 2.3 compute the corresponding parameters of the CE framework. 
However, it is unclear whether the residuals $\boldsymbol{Z}_{\mid i,s}$ follow a multivariate normal distribution, and if so, for which parameters. 
Since the proposed method relies on this assumption, it is crucial to comment on the validity of the simulation setting and the implications for the method’s performance.}

\noindent \textcolor{teal}{\bf Response}: \\

%%% Item %%% Item %%% Item %%% Item %%% Item %%% Item %%% Item %%%

\noindent \textcolor{teal}{
2. In Section 3.2, it is unclear whether the introduced model fits within the CEV framework.
Additionally, what are the parameters $\alpha_{\mid i, s}$ and $\beta{_\mid i, s}$ for $i \in \{1, \ldots, d\}$ and $s \in \{1, \ldots, D\}$?
}

\noindent \textcolor{teal}{\bf Response}: \\

%%% Item %%% Item %%% Item %%% Item %%% Item %%% Item %%% Item %%%

\todo{Fix references}
\noindent \textcolor{teal}{
3. In Section 3.2.1, it is unclear why the comparison between the proposed method and that of
Vignotto, Engelke, and Zscheischler 2021 adds significant value to the paper. 
While the method operates within the Asymptotic Dependence (AD) framework, it is not formulated within the CEV framework. 
Perhaps, an interesting comparison could involve two simulation settings: one under Asymptotic Independence (AI), i.e., $-1 < \alpha_{j \mid i, s} < 1$ for all $i, j \in \{1, \ldots, d\}$ with $j \ne i$ and $s \in \{1, \ldots, D\}$, and another under Asymptotic Dependence (AD), i.e. $\alpha_{j \mid i, s} = 1$ for all $i, j \in \{1, \ldots, d\}$ with $j \ne i$ and $s \in \{1, \ldots, D\}$. 
One could then compare the performance of the two methods in these distinct settings. 
Since the proposed method encompasses both AI and AD, while Vignotto, Engelke, and Zscheischler 2021 only handles AD, we would expect the proposed method to perform well in both settings, whereas the competitor would only perform well in the AD scenario. 
It is possible that the results in Section 3.2.1 already include such comparisons, but as presented, these findings are not clearly discussed.
}

\noindent \textcolor{teal}{\bf Response}: \\

%%% Item %%% Item %%% Item %%% Item %%% Item %%% Item %%% Item %%%

\noindent \textcolor{teal}{
4. In Section 3.2.2, it is unclear what the added value of this setting is. 
As the authors note,
since $\rho^{s}_t$ is the same across all variable pairs, increasing the dimension d effectively increases the amount of available data. 
Consequently, the results presented in this section may simply reflect the effect of having more data, rather than a fundamental property of the proposed method. 
As suggested by the authors, it would be more informative to consider a setting where 
$\rho^{s}_t$ differs across variable pairs. 
I strongly recommend that the authors explore this direction. 
This would provide stronger evidence that ``the proposed method scales effectively to higher dimensions and performs well in multivariate settings.'' 
Otherwise, the conclusion should be interpreted with caution or possibly omitted, as its current basis is not robust.
}

\noindent \textcolor{teal}{\bf Response}: \\

%%% Item %%% Item %%% Item %%% Item %%% Item %%% Item %%% Item %%%

\noindent \textcolor{teal}{
My third concern pertains to the data used in the application and the lack of rigor in assessing
whether these data fully comply with the modeling assumptions introduced in Section 2, an issue
that is not addressed by the authors.
}

\noindent \textcolor{teal}{
1. As far as I understand, the data are collected from two different sources: precipitation from
Met Éireann and ERA5 reanalysis. 
Presumably, the former are observational data, while the latter are modeled data. 
The use of precipitation data from Met Eireann and wind data from ERA5 raises important methodological questions regarding the potential introduction of biases in the dependence structure between these variables. 
Observational and reanalysis data differ in spatial/temporal resolution, measurement uncertainty, and processing methods, which may artificially inflate or diminish apparent correlations. 
The authors should justify this choice, discuss its impact on the results, and consider either using a single data source or applying harmonization techniques to mitigate these biases.
}

\noindent \textcolor{teal}{\bf Response}: \\

%%% Item %%% Item %%% Item %%% Item %%% Item %%% Item %%% Item %%%

\noindent \textcolor{teal}{
2. The spatial resolution of ERA5 may be relatively coarse for the Island of Ireland. 
For example, at locations 1, 2, 3, and 5 in the right panel of Figure 6, there may be multiple points contained within the same $0.25 \times 0.25$ degree pixel, potentially resulting in identical time series for wind speed at these locations. 
Does this actually occur, or were the selected locations chosen carefully to avoid such overlap? 
If not, would the conclusions of the study remain unaffected?
}

\noindent \textcolor{teal}{\bf Response}: \\

%%% Item %%% Item %%% Item %%% Item %%% Item %%% Item %%% Item %%%

\noindent \textcolor{teal}{
3. The proposed method is developed within the CEV framework defined in Equation (5), where
$\boldsymbol{Z}_{\mid i, s}$ is assumed to follow a multivariate Gaussian distribution with mean vector $\mu_{\mid i, s} \in \mathbb{R}$ and covariance matrix 
$\Sigma_{\mid i, s}$ for $i = 1, 2$ and $s = 1, \ldots, 59$. 
Could the authors provide, in Section 4.2, more quantitative justification supporting the assumption that
\[
  \boldsymbol{Z}_{\mid i, s} = \frac{\boldsymbol{Y}_{-i}-\hat{\alpha}_{\mid i, s}\boldsymbol{Y}_i}{\boldsymbol{Y}_{i}^{\hat{\beta}_{i, s}}}
\]
for $\boldsymbol{Y}_i > u_{i, s}$ can be reasonably modeled by a multivariate normal distribution? 
Indeed, the above expression is an estimator of the distribution of $G_{\mid i, s}$, as stated in Keef, Papastathopoulos, and Tawn 2013, Section 2.4. 
Given the large number of variables considered, the authors could, for instance, provide one or two representative Q–Q plots comparing the fitted Gaussian model to the empirical observations, or diagnostic plots of the residuals. 
The results for all 118 cases could then be summarized using a boxplot of the 118 coefficients of determination between the fitted values and the observed values. 
If I am not mistaken, a similar diagnostic analysis has been carried out in Ross et al. 2018, Figure 6 in a same setting to that considered in the present paper, suggesting that such an assessment should be feasible in practice.
}

\noindent \textcolor{teal}{\bf Response}: \\


%%% Item %%% Item %%% Item %%% Item %%% Item %%% Item %%% Item %%%

\noindent \textcolor{teal}{
4. The authors provide a thorough spatial analysis of the clusters, which helps to assess their
spatial coherence. 
However, is it also possible to interpret the clusters in terms of their dissimilarity? 
For instance, the second cluster (shown in blue) in Figure 9 appears to consist of stations with very small dissimilarity values. 
Can this observation be interpreted within the CEV framework?
}

\noindent \textcolor{teal}{\bf Response}: \\

%%% Section %%% Section %%% Section %%%
%%% Section %%% Section %%% Section %%%
%%% Section %%% Section %%% Section %%%
\section*{{Minor Comments}}

%%% Item %%% Item %%% Item %%% Item %%% Item %%% Item %%% Item %%%
{\color{teal}
\noindent 
On page 11, $h(y \mid \boldsymbol{Y}_{i,s} > u_i)$ is defined as the conditional density of $\boldsymbol{Y}_{i,s} \mid (Y_{i,s} > u_i)$. 
However, it is unclear what is meant by $h(y \mid \boldsymbol{Y}_{i,s^*} > u_i)$. 
Is this intended to be the conditional density of $\boldsymbol{Y}_{i,s} \mid (\boldsymbol{Y}_{i,s^*} > u_i)$, or of $\boldsymbol{Y}_{i,s^*} \mid (\boldsymbol{Y}_{i,s^*} > u_i)$? 
If the first interpretation is intended, it would be helpful to provide more details or justification for why $h(y \mid \boldsymbol{Y}_{i,s} > u_i) = h(y \mid \boldsymbol{Y}_{i,s^*} > u_i)$.
}

\noindent \textcolor{teal}{\bf Response}: \\

%%% Item %%% Item %%% Item %%% Item %%% Item %%% Item %%% Item %%%
\todo{Check that eJSGa is correct here}
{\color{teal}
\noindent 
  In Section 2.4, Algorithm 1 is described as depending on the conditioning variable $i$ as an input. 
However, in Section 3.1 it is applied using the matrix
\[
  M = \frac{1}{d}\sum_{i=1}^d M^{(i)}.
\]
In this context, it is unclear how $\eJSGa{i,s,s^*}$ for $i \in \{1, \ldots, d\}$ and $s, s^* \in \{1, \ldots, D\}$ is defined, since the divergence depends on i, but $M$ does not. 
Presumably, the authors intend to refer to
\[
  \frac{1}{d}\sum^d_{i=1}\eJSGa{i, s, s^*}
\]
for two specific sites $s, s^* \in \{1, \ldots, D\}$. 
As currently stated, Algorithm 1 does not flexibly accommodate this situation. 
A slight rewriting of the algorithm would make it more general, allowing it to handle both the per-$i$ and aggregated cases.
}

\noindent \textcolor{teal}{\bf Response}: \\

%%% Item %%% Item %%% Item %%% Item %%% Item %%% Item %%% Item %%%

{\color{teal}
\noindent
In Section 3.1, it is unclear what the term ``pre-clustering'' refers to. 
My understanding is that no clustering method has been applied. 
If that is the case, why does the ``pre-clustering'' panel in Figure 1 display three boxplots for each dependence quantile? 
Could the authors clarify what is being shown in this panel?
}

\noindent \textcolor{teal}{\bf Response}: \\

%%% Item %%% Item %%% Item %%% Item %%% Item %%% Item %%% Item %%%

{\color{teal}
\noindent
 On page 21, could the authors be more precise about where Bengtsson, Hodges, and Keenlyside
2009 indicates that the weekly temporal scale reflects the lag between extreme precipitation
and wind speed during storm events?
}

\noindent \textcolor{teal}{\bf Response}: \\

%%% Item %%% Item %%% Item %%% Item %%% Item %%% Item %%% Item %%%

{\color{teal}
\noindent
In Figure 7, the estimates of $\alpha_{j \mid i, s}$ for all $i, j \in \{1, 2\}$ with $j \ne i$ lie between 0 and 1, which implies asymptotic independence between precipitation and wind speed (see, e.g., Keef, Papastathopoulos, and Tawn 2013, Section 2.3). 
Is this not in contradiction with the statement in Section 4.1 that all sites exhibit positive extremal dependence between precipitation and wind speed, or am I missing an important distinction?
}

\noindent \textcolor{teal}{\bf Response}: \\

%%% Item %%% Item %%% Item %%% Item %%% Item %%% Item %%% Item %%%

{\color{teal}
\noindent
 In Figure 8, I think the last item should be ``$> 0.08$'' instead of ``$> 0.8$''.
}

\noindent \textcolor{teal}{\bf Response}: \\

\clearpage
\newpage


%%%%%%%%%%%%%%%%%%%%%%%%%%%%%%%%%%%%%%%%%%%%%%%%
%%% Reviewer 2 %%% Reviewer 2 %%% Reviewer 2 %%%
%%%%%%%%%%%%%%%%%%%%%%%%%%%%%%%%%%%%%%%%%%%%%%%%
\noindent {\Large {\bf Responses to Reviewer 2}}\\

\noindent \textcolor{blue}{
This manuscript addresses the challenge of interpreting and comparing extremal dependence structures across high-dimensional multivariate data, particularly in environmental contexts.
The authors build on the conditional extremes (CE) framework, which models joint tail behavior, and introduce a novel clustering method that employs a closed-form skew-geometric Jensen-Shannon divergence to quantify dissimilarities between CE models across sites or variables.
Their approach outperforms existing bivariate clustering methods in simulations and extends naturally to higher dimensions. 
Applying this to Irish meteorological data, the authors identify spatially coherent regions with similar extremal dependence between precipitation and wind speeds, providing new insights into the spatial structure of environmental extremes.
}

\noindent \textcolor{blue}{
I found the manuscript interesting, well written, and well structured. 
The proposed clustering method clearly addresses an important gap in multivariate extreme clustering. 
From a methodological perspective, the contribution appears to rely more on a thoughtful combination of existing concepts—namely, the CE framework and the skew-geometric Jensen-Shannon divergence—than on fundamentally new methodological developments. 
Nevertheless, the proposed approach, together with the synthetic and real case studies, yields interesting results and represents a worthwhile contribution.
}

\noindent \textcolor{blue}{\bf Response}: Thank you for these positive comments.\\

%%% Section %%% Section %%% Section %%%
%%% Section %%% Section %%% Section %%%
%%% Section %%% Section %%% Section %%%

\section*{{Minor Comments}}

%%% Item %%% Item %%% Item %%% Item %%% Item %%% Item %%% Item %%%

{\color{blue}
\noindent
1. Some parts of the Introduction could benefit from improved flow. In particular, on p.3 (second paragraph), the statement ``we have limited tools for identifying random vectors\ldots'' is made without immediately providing examples (which appear later in the Introduction). 
In addition, the paragraph concludes by introducing the use of the CE framework as a contribution, but without sufficient justification at that stage, which is only partly developed later.
}

\noindent \textcolor{blue}{\bf Response}: \\

%%% Item %%% Item %%% Item %%% Item %%% Item %%% Item %%% Item %%%

{\color{blue}
\noindent
2. On p.9, at the end of the first paragraph, I believe this should be specified for a given $i, 1 \le i \le d$. 
}

\noindent \textcolor{blue}{\bf Response}: \\

%%% Item %%% Item %%% Item %%% Item %%% Item %%% Item %%% Item %%%

{\color{blue}
\noindent
3. On p.9, the last sentence (``Its computational efficiency \ldots'') would be more appropriately placed in the Introduction.
}

\noindent \textcolor{blue}{\bf Response}: \\

%%% Item %%% Item %%% Item %%% Item %%% Item %%% Item %%% Item %%%

{\color{blue}
\noindent
4. On p.11, since Monte Carlo methods are required to estimate Eq.(11), providing an estimate of the associated computational cost would be useful.
}

\noindent \textcolor{blue}{\bf Response}: \\

%%% Item %%% Item %%% Item %%% Item %%% Item %%% Item %%% Item %%%

{\color{blue}
\noindent
5. On p.13, Algorithm 1 uses $M^(i)$ as input, but $M^(i)$ is never referenced in the calculations inside the algorithm; I believe this is a mistake.
}

\noindent \textcolor{blue}{\bf Response}: \\

%%% Item %%% Item %%% Item %%% Item %%% Item %%% Item %%% Item %%%

{\color{blue}
\noindent
6. On p.14, fourth line: I believe it should be the correlation matrix rather than the covariance matrix, since the off-diagonal entries are pairwise correlations.
}

%%% Item %%% Item %%% Item %%% Item %%% Item %%% Item %%% Item %%%

{\color{blue}
\noindent
7. End of p.14: I suggest: The explanations of the pre-clustering and post-clustering steps should be clearer and more straightforward.
}

\noindent \textcolor{blue}{\bf Response}: \\

%%% Item %%% Item %%% Item %%% Item %%% Item %%% Item %%% Item %%%

{\color{blue}
\noindent
8. On p.15, sec. 3.2: I suggest: I would like a more detailed justification for this specific mixture choice - for example, why not use a mixture of Student's t distributions only?
}

\noindent \textcolor{blue}{\bf Response}: \\

%%% Item %%% Item %%% Item %%% Item %%% Item %%% Item %%% Item %%%

{\color{blue}
\noindent
9. On p.20, sec. 4.1: could you compare your results with those of Vignotto et al. (2021)?
}

\noindent \textcolor{blue}{\bf Response}: \\

\bibliographystyle{apalike}
{\small
\bibliography{ref}
}
\end{document}
